\documentclass[11pt,a4paper]{article}
\usepackage{fontspec,amsmath,amssymb,amsthm,microtype}
\setmainfont{lmroman10-regular.otf}[BoldFont=lmroman10-bold.otf,ItalicFont=lmroman10-italic.otf,BoldItalicFont=lmroman10-bolditalic.otf]
\usepackage[margin=25mm]{geometry}
\usepackage[unicode,hidelinks]{hyperref}
\hypersetup{pdfauthor={LangevinLean Project}}
\setlength{\emergencystretch}{3em}
\setlength{\parskip}{0.4em}
\setlength{\parindent}{0pt}
\allowdisplaybreaks[2]
\raggedbottom
\title{\texorpdfstring{Exact kinetic Gibbs stationarity from the initial/\allowbreak{}noise laws}{Exact kinetic Gibbs stationarity from the initial/noise laws}}
\author{LangevinLean Project}
\date{}
\begin{document}\maketitle
Draft mathematical exposition; these notes may contain mistakes.


The stationarity argument starts with a continuous L\ensuremath{{}^{2}} kinetic integral solution whose joint input law is the product of the target initial phase law and the complete canonical Ornstein–Uhlenbeck innovation stream. The independent product construction supplies this joint law. Convergence of splitting approximations, vanishing stationary weak defects and equality of the limiting endpoint laws then prove invariance.



The splitting approximants have a vanishing target weak discrepancy and converge to the constructed kinetic solution. Their common limit identifies the solution’s full Gibbs phase law at every time.


\section{\texorpdfstring{Derive exact stationarity through the splitting limit}{Derive exact stationarity through the splitting limit}}


Let \(\pi\) be the normalized Gaussian-times-Gibbs target, and let \(Y\) solve the kinetic integral equation driven by the continuous L\ensuremath{{}^{2}} input \(B\). On a mesh \(h_n=T/(n+1)\), form the concrete Verlet–OU approximation \(S_n\) using the filtered increments of this same \(B\). The proved Banach splitting theorem gives \(S_n\to Y_T\) in the product of L\ensuremath{{}^{2}} spaces.



The full initial/\allowbreak{}innovation joint law is the explicit product \(\pi\otimes Q_{h_n}^{\otimes\mathbb N}\). The source-law transfer theorem, together with the almost-everywhere vector-chain representation, identifies the law of \(S_n\) with the defined Markov iterate. Its bounded-test error is at most \[2(n+1)h_n^2 K A=\frac{2T^2KA}{n+1},\qquad K=\frac{M(2M+11)}8,\quad A=\int(\|v\|^2+\|\nabla f(x)\|^2)\,d\pi.\] The moment is finite by the proved target tail theorem. The mesh is eventually at most one, so every quadratic rejection estimate applies, and the displayed scalar error tends to zero.



The phase-law continuity theorem now identifies the L\ensuremath{{}^{2}} limiting phase law with \(\pi\). The argument constructs the true approximation sequence, applies the proved convergence theorem, rewrites its chain law using the explicit initial/\allowbreak{}noise product distribution, rewrites the bounded test expectation through the pushforward, invokes the proved Markov weak-error bound, and checks the scalar \(T^2/(n+1)\) decay. It then applies dual-characteristic uniqueness for finite measures. No stationary solution, local-error conclusion, output-law independence field, or numerical-limit conclusion is assumed. The remaining input-law hypothesis is a concrete source probability identity, supplied by the canonical independent-product construction.



The constructed splitting chain has the exact independent thermal-noise law and its endpoint converges in second mean to the kinetic solution. Its bounded-test target discrepancy vanishes along the same refinement. Characteristic functions therefore identify the limiting phase law with the Gibbs target, proving stationarity of the continuous solution.


\end{document}
